\section{Introduction}
LLM agents that operate over long horizons are routinely equipped with memory modules: vector stores, summarization pipelines, graph memories, and retrieval-augmented generation (RAG) stacks designed to surface the right past fact at question time~\cite{locomo2024,longmemeval2024,amem2025}. These systems primarily target \emph{retrospective} memory---optimizing retrieval similarity $\mathrm{sim}(e,q)$ to answer ``what happened?'' from chat history. In everyday assistance, however, users also issue deferred instructions: take medicine at 21:00, book a slot when the portal opens, reply once an email arrives, cancel the reminder if plans change. Success here is not measured by retrieving a needle from a haystack, but by \emph{acting at the right future moment} while still advancing concurrent tasks, and by \emph{not} acting when the intention has been canceled or the cue is only a lure.

Cognitive science has long distinguished retrospective recall from \emph{prospective memory} (PM)---remembering to execute delayed intentions under cue- or time-based triggers amid ongoing activity~\cite{einstein2005prospective,rendell2000virtual}. PM-Bench~\cite{pmbench2026} systematizes this gap for LLM agents via a seven-day Virtual Week: at each step the agent must choose a mandatory ongoing activity and optionally fire a set of prospective actions against a ground-truth \emph{due set} $D_t$, while hidden state channels (clock, email, portals, \ldots) are invisible unless actively queried. Empirically, frontier models remain fragile: channel-required hit rates are near floor, cross-day and update-sensitive intentions are frequently missed, and scaffolds that spam monitoring queries inflate false positives and degrade Set-F1. Among eight published scaffolds, optional heartbeat currently achieves the best macro Set-F1 ($65.1\%$); auto-heartbeat raises update-sensitive hits but floods false positives, and hierarchical union-query issues thousands of channel reads yet still misses channel-required dues~\cite{pmbench2026}. No universal off-the-shelf pattern---plain single-agent prompting, todo ledgers, periodic heartbeats, or hierarchical union-query---dominates every diagnostic slice.

We contend that the remaining gap is architectural as well as evaluative. Retrospective memory treats past events as content to be retrieved under a query; prospective competence requires \emph{conditional} records that remain dormant until a trigger fires, then transition through a lifecycle that supports modification and cancellation. Conflating the two---storing intentions as ordinary episodic text and hoping RAG or long context will ``notice'' when they are due---fails exactly where PM-Bench isolates difficulty: monitoring under uncertainty, discriminating lures, and preserving commitments across days and updates. Our contribution is not to rediscover that agents need prospective memory---PM-Bench already established that---but to \textbf{instantiate typed conditional memory} as a reproducible scaffold on this benchmark.

\begin{figure*}[t!]
\centering
\includegraphics[width=0.9\textwidth]{figure/promem_overview.pdf}
\caption{Overview of ProMem. Deferred intentions are stored as conditional memory triples $(\varphi,\alpha,\sigma)$ in a Prospective Intention Store (PIS), separate from an episodic retrospective bank. A lifecycle manager updates $\sigma$ on announce, cue, fire, complete, cancel, and modify events. A proactive trigger monitor decides when to query hidden channels; a due-set scorer then selects executable actions (including lure rejection) while the agent still chooses an ongoing activity each step, following the PM-Bench query-then-act protocol.}
\label{fig:promem-overview}
\end{figure*}

To that end we introduce \textbf{conditional memory} as a first-class abstraction---structured intentions $I=(\varphi,\alpha,\sigma,\mathrm{meta})$ whose success criterion is trigger satisfaction $\mathrm{sat}(\varphi,e_t)$, not query--document similarity---and instantiate it with \textbf{ProMem} (Figure~\ref{fig:promem-overview}): dual memory stores, an intention lifecycle manager, a proactive trigger-monitoring policy, and a due-set scorer integrated with ongoing activity choice. Our contributions are:
\begin{itemize}
    \item We formalize \emph{conditional memory} as an abstraction orthogonal to retrospective episodic banks, with due sets $D_t$ and Set-F1 as the natural objective, anchored on PM-Bench's published failure modes (channel monitoring, cross-day carry-over, update sensitivity).
    \item We propose ProMem: dual stores (episodic bank + Prospective Intention Store), explicit lifecycle transitions, an intention-conditioned monitor $\pi_{\mathrm{mon}}$ for hidden channels, and a lure-aware due-set scorer $\pi_{\mathrm{due}}$---targeting the FP explosion and channel under-hit that PM-Bench reports under todo ledgers and fixed-interval heartbeats.
    \item We outline a PM-Bench-centric evaluation protocol with baselines matching all eight published scaffolds, ablations isolating each ProMem component, and slice metrics (channel / cross-day / update-sensitive); numeric results are pending final runs and we do not claim dominance over optional heartbeat without measured Set-F1.
\end{itemize}
